\documentclass[11pt]{article}
\usepackage{amsmath,amssymb}
\title{P1\_N6\_v2 series: scope correction (Room 18, confirmed by Reviewer BB)}
\date{2026-09-27}
\begin{document}
\maketitle

\section*{Correction to the whole P1\_N6\_v2 series (Rooms 8, 11, 13, 14, 16, 17)}

\textbf{Accumulation of poles at $\zeta_0=e(1/6)$ was never open.} \texttt{paper2a.tex} itself
already states, verbatim, at the point where the N=6 obstruction is first discussed (the
\texttt{rem:arcobstruction} paragraph following $N=6<16$): ``Accumulation of the zeros of $B$ at
$e^{2\pi i/6}$, and of the poles of $U,V$ at $\pm\sqrt{e^{\pm2\pi i/6}}$, still holds
unconditionally-in-method by closure from Theorem~\texttt{thm:arc}/Corollary~\texttt{cor:arcUV}
... what is open at $1/6$ is only the direct, single-point, tie-ray statement.'' This is a
straightforward instance of the closed-derived-set argument: $1/6\approx0.1667$ lies strictly
inside the proved arc $[0.143,0.857]$ (\texttt{thm:arc}) and at the closed endpoint of
$[1/6,5/6]$ (\texttt{cor:arcUV}), both proved dense-hence-closed accumulation sets. The same fact
is also stated in \texttt{P1i\_lemma.tex} (lines 27, 186) and in the original
\texttt{P1\_N6/room1\_seat\_erdos.tex} (line 20).

\textbf{None of Rooms 8, 11, 13, 14, 16, 17 cited this}, and their stated goals (``close N=6,''
``N=6 remains OPEN'') conflated the already-proved accumulation statement with a strictly
\emph{stronger}, genuinely open question: \textbf{the direct tie-ray statement at N=6 itself} --
an explicit, N=6-specific sequence of zeros along the tie ray, established via an N=6 analogue of
\texttt{prop:qi}'s own contour/Rouch\'e construction, rather than via the closure argument.

\textbf{Nothing mathematical in the series is wrong.} Room 17's separation-constant bound
($D(a',\varphi)\ge1-e^{-a'}$, band-uniform, confirmed by 3 independent reviewers) is a real,
correct fact about a real sub-piece of the (still hypothetical) direct construction. What needs
correcting is only the framing: this work is progress toward the \emph{direct tie-ray statement},
not toward accumulation (which needs nothing further).

\textbf{Secondary correction (Room 18, confirmed):} Room 11's stated verdict --- ``the ray lemma
gives no lower bound on the separation constant anywhere along the resonant legs'' --- was an
overgeneralization. Room 11's own detailed argument only established this for the \emph{initial
stretch} near the starting point $x_c$; Room 17's bound shows $\delta\ge1-e^{-2\operatorname{Re}(x)}$
holds cleanly on the entire $\operatorname{Re}(x)>0$ portion of both legs. \textbf{The true
remaining gap for the hypothetical direct construction is narrower still: specifically the
$\operatorname{Re}(x)\lesssim0$ crossing region}, i.e.\ an N=6 analogue of
\texttt{lem:qi-bounds}(ii), which nobody has attempted.

\textbf{No paper2a edit needed} -- its own wording was already correct; only this repository's
internal room notes needed the correction, recorded here.

\section*{Corrected status, going forward}

\begin{itemize}
\item Accumulation at $e(1/6)$ (and $\pm\sqrt{\cdot}$ for U,V): \textbf{PROVED}, always was, via
closure. Nothing to do here.
\item Direct tie-ray statement at N=6: \textbf{OPEN}. Progress so far: the separation-constant
piece on $\operatorname{Re}(x)>0$ is solved (Room 17). The $\operatorname{Re}(x)\lesssim0$
crossing-region estimate (an N=6 analogue of \texttt{lem:qi-bounds}(ii)) is the next concrete,
well-scoped target, not yet attempted.
\end{itemize}

\end{document}
